\chapter{Chapter05 — Berry Esseen Rescaling}
%%% generated by the blueprint pipeline from FABL.Chapter05.BerryEsseenRescaling
\section{Overview}

Declarations of the Lean module \texttt{FABL.Chapter05.BerryEsseenRescaling}.

\section{Declarations}

\begin{theorem}[Berry–Esseen bound with nonzero means]
\lean{FABL.exercise5_17}
\proofsource{Analysis-of-Boolean-Functions-by-Ryan-ODonnell}{pdf-ffbe9c7b3306-p119-b002}
\label{FABL.exercise5_17}
\leanok
\difficulty{6}
Let $\{X_i\}_{i\in\iota}$ be a finite family of independent, measurable real random
variables on a probability space, each with finite second moment. Suppose the sum of
their variances is positive and each centered variable $X_i-\mathbb{E}X_i$ has finite
third absolute moment. For every $x\in\mathbb{R}$, the absolute difference between the
cumulative distribution function of $\sum_i X_i$ at $x$ and that of the Gaussian with
mean $\sum_i\mathbb{E}X_i$ and variance $\sum_i\operatorname{Var}(X_i)$ is at most
$C\,\sum_i\mathbb{E}|X_i-\mathbb{E}X_i|^3/(\sum_i\operatorname{Var}(X_i))^{3/2}$, where
$C$ is the third-moment Berry–Esseen constant.
\end{theorem}
